\chapter{Generalized symmetries and anomaly inflow}\label{ch:gensym}

Chapter~\ref{ch:dw} gave us a clean discrete model of topological response and anomaly inflow. The purpose of the present chapter is to extract the general language behind that example. Rather than starting from differential geometry or abstract category theory, we begin from a fact every quantum-field-theory student already knows: an ordinary global symmetry has a conserved current and an associated conserved charge. The key step is to rewrite that familiar statement in the language of topological defects.

\medskip

\noindent Once phrased that way, the jump from ordinary symmetries to higher-form symmetries is almost not a jump at all. One simply changes the dimension of the charged operator and of the topological defect that measures it. This chapter follows the defect-language viewpoint emphasized in the modern generalized-symmetry literature \cite{GKSW2015,Bhardwaj2023Lectures}. Section~\ref{sec:zero-form-defects} recasts ordinary $0$-form symmetries in that language; Section~8.2 will then generalize the same pattern to arbitrary $q$-form symmetries. Later sections will return to Chapter~\ref{ch:dw} and reinterpret its twisted finite-group theories as concrete anomaly-inflow laboratories.

\section{Ordinary \texorpdfstring{$0$}{0}-form symmetries in defect language}\label{sec:zero-form-defects}

The slogan of this section is simple:
\begin{quote}
an ordinary global symmetry can be represented by a topological operator supported on a codimension-$1$ manifold.
\end{quote}
That statement is often taken as the definition of a generalized symmetry. For ordinary symmetries, however, it is not a new postulate; it is just Noether's theorem rewritten in a form that generalizes.

\subsection{From a conserved current to a topological surface operator}

\begin{assumption}[Standing assumptions for Section~\ref{sec:zero-form-defects}]\label{ass:zero-form-standing}
Let $X$ be an oriented $d$-dimensional spacetime manifold, and suppose the quantum field theory on $X$ has a continuous internal $0$-form symmetry with conserved Noether current $j^\mu$. We work in Euclidean signature when discussing defect operators, so the conservation law is written equivalently as
\begin{equation}\label{eq:noether-conservation-forms}
\partial_\mu j^\mu = 0
\qquad\Longleftrightarrow\qquad
d(\star j)=0.
\end{equation}
For notational simplicity we first discuss a $U(1)$ symmetry, or more generally a single one-parameter subgroup of a continuous symmetry group.
\end{assumption}

\begin{definition}[Charge carried by a codimension-$1$ surface]\label{def:zero-form-charge}
Let $\Sigma \subset X$ be a closed oriented codimension-$1$ submanifold, so $\dim \Sigma = d-1$. The Noether charge associated with $\Sigma$ is
\begin{equation}\label{eq:QSigma}
Q_\Sigma := \int_\Sigma \star j.
\end{equation}
If $\Sigma$ is a constant-time spatial slice in Lorentzian signature, then \eqref{eq:QSigma} is the usual conserved symmetry charge.
\end{definition}

\begin{definition}[Symmetry defect operator]\label{def:zero-form-defect}
For a group element in the chosen one-parameter subgroup, written as
\begin{equation*}
g=e^{i\alpha},
\end{equation*}
define the corresponding codimension-$1$ symmetry operator by
\begin{equation}\label{eq:UgSigma}
U_g(\Sigma) := \exp\!\bigl(i\alpha Q_\Sigma\bigr).
\end{equation}
We will refer to $U_g(\Sigma)$ as the \emph{symmetry defect} or \emph{topological symmetry operator} supported on $\Sigma$.
\end{definition}

\begin{remark}[Why codimension $1$ is the right dimension]\label{rem:codim-one-zero-form}
The $(d-1)$-form $\star j$ can be integrated only over $(d-1)$-dimensional submanifolds. The operator $U_g(\Sigma)$ therefore lives on a codimension-$1$ support. In the generalized-symmetry dictionary, that is exactly what one expects for a $0$-form symmetry: the charged operators are pointlike, and the symmetry operator linking them is one codimension higher.
\end{remark}

\begin{proposition}[Current conservation implies topological invariance]\label{prop:zero-form-topological}
Let $\Sigma_0$ and $\Sigma_1$ be closed oriented codimension-$1$ submanifolds that are homologous inside a region $V \subset X$ containing no charged operator insertions. Then
\begin{equation}\label{eq:QSigma-invariant}
Q_{\Sigma_0}=Q_{\Sigma_1},
\qquad
U_g(\Sigma_0)=U_g(\Sigma_1).
\end{equation}
Equivalently, the symmetry operator $U_g(\Sigma)$ may be deformed continuously without changing correlation functions, provided the deformation does not cross a charged insertion.
\end{proposition}

\begin{proof}
Because $\Sigma_0$ and $\Sigma_1$ are homologous inside $V$, there exists an oriented region $W \subset V$ with boundary
\begin{equation*}
\partial W = \overline{\Sigma_0}\sqcup \Sigma_1.
\end{equation*}
Applying Stokes' theorem and using \eqref{eq:noether-conservation-forms}, we obtain
\begin{align*}
Q_{\Sigma_1}-Q_{\Sigma_0}
&=
\int_{\Sigma_1}\star j-\int_{\Sigma_0}\star j \\
&=
\int_{\partial W}\star j \\
&=
\int_W d(\star j) \\
&= 0.
\end{align*}
Therefore $Q_{\Sigma_1}=Q_{\Sigma_0}$. Exponentiating both sides yields
\begin{equation*}
U_g(\Sigma_1)=\exp(i\alpha Q_{\Sigma_1})
=
\exp(i\alpha Q_{\Sigma_0})
=
U_g(\Sigma_0).
\end{equation*}
This is precisely the topological-invariance statement.
\end{proof}

\begin{remark}[What is and is not being used]\label{rem:zero-form-ingredients}
Proposition~\ref{prop:zero-form-topological} uses only the conservation law and Stokes' theorem. No metric information survives in the deformation statement except what was needed to define the Hodge star in the first place. This is why the defect-language reformulation is robust: once the topological operator has been identified, the later generalized-symmetry formalism no longer depends on the details of the original Noether-current derivation.
\end{remark}

\subsection{Charged local operators and the crossing relation}

\begin{definition}[Local operator of charge \texorpdfstring{$q$}{q}]\label{def:charged-local-operator}
A local operator $\mathcal O_q(x)$ inserted at a point $x\in X$ has charge $q$ if, for a sufficiently small sphere $S_x^{d-1}$ surrounding $x$,
\begin{equation}\label{eq:local-charge}
Q_{S_x^{d-1}}\,\mathcal O_q(x)=q\,\mathcal O_q(x).
\end{equation}
Equivalently, in canonical language the conserved charge acts by
\begin{equation*}
[Q,\mathcal O_q(x)] = q\,\mathcal O_q(x).
\end{equation*}
\end{definition}

\noindent The point of the defect language is that the symmetry action can now be read geometrically: a codimension-$1$ symmetry defect acts on a point operator by whether or not the defect surface links that point.

\begin{proposition}[Crossing a charged local operator produces the symmetry action]\label{prop:zero-form-crossing}
Let $\Sigma_{<}$ and $\Sigma_{>}$ be two codimension-$1$ surfaces that differ by moving the defect across a single insertion of $\mathcal O_q(x)$. Then
\begin{equation}\label{eq:crossing-relation}
U_g(\Sigma_{>})\,\mathcal O_q(x)
=
e^{i\alpha q}\,\mathcal O_q(x)\,U_g(\Sigma_{<}).
\end{equation}
In other words, the topological defect acts on a charged local operator by the expected group action.
\end{proposition}

\begin{proof}
Choose a small sphere $S_x^{d-1}$ around the insertion point $x$ and a region $W$ whose boundary is
\begin{equation*}
\partial W = \overline{\Sigma_{<}}\sqcup \Sigma_{>}\sqcup \overline{S_x^{d-1}}.
\end{equation*}
The extra boundary component appears because the deformation from $\Sigma_{<}$ to $\Sigma_{>}$ passes across the point $x$, so the region must be punctured there. Applying Stokes' theorem away from the insertion gives
\begin{align*}
0
&=
\int_W d(\star j) \\
&=
\int_{\Sigma_{>}}\star j
-\int_{\Sigma_{<}}\star j
-\int_{S_x^{d-1}}\star j.
\end{align*}
Hence
\begin{equation}\label{eq:charge-crossing}
Q_{\Sigma_{>}} = Q_{\Sigma_{<}} + Q_{S_x^{d-1}}.
\end{equation}
Now let this identity act on $\mathcal O_q(x)$. By Definition~\ref{def:charged-local-operator},
\begin{equation*}
Q_{S_x^{d-1}}\,\mathcal O_q(x)=q\,\mathcal O_q(x),
\end{equation*}
so exponentiating \eqref{eq:charge-crossing} yields
\begin{align*}
U_g(\Sigma_{>})\,\mathcal O_q(x)
&=
\exp\!\bigl(i\alpha Q_{\Sigma_{>}}\bigr)\,\mathcal O_q(x) \\
&=
\exp\!\bigl(i\alpha Q_{\Sigma_{<}}\bigr)\exp\!\bigl(i\alpha Q_{S_x^{d-1}}\bigr)\,\mathcal O_q(x) \\
&=
\exp(i\alpha q)\,\mathcal O_q(x)\,\exp\!\bigl(i\alpha Q_{\Sigma_{<}}\bigr) \\
&=
e^{i\alpha q}\,\mathcal O_q(x)\,U_g(\Sigma_{<}),
\end{align*}
which is \eqref{eq:crossing-relation}.
\end{proof}

\begin{remark}[Why this is called a \texorpdfstring{$0$}{0}-form symmetry]\label{rem:why-zero-form}
The charged operator in Proposition~\ref{prop:zero-form-crossing} is supported at a point, so it is a $0$-dimensional operator in spacetime. This is the meaning of the label ``$0$-form.'' The associated symmetry operator is supported on codimension $q+1=1$. In a $d$-dimensional spacetime, its support therefore has dimension
\begin{equation*}
d-(q+1)=d-1.
\end{equation*}
Section~8.2 will keep exactly this codimension rule and change only the dimension of the charged operator.
\end{remark}

\begin{remark}[Why the defect viewpoint is more general than Noether]\label{rem:noether-vs-defect}
For continuous symmetries, Noether's theorem gives a conserved current and therefore a direct route to the topological operator $U_g(\Sigma)$. For discrete symmetries there need not be any local current at all, but the codimension-$1$ topological defect still makes sense. This is why the defect language is the right starting point for generalized symmetries: it contains the ordinary Noether story as a special case, but it is not limited to it.
\end{remark}

\begin{remark}[Forward pointer]\label{rem:zero-form-forward}
Section~\ref{sec:zero-form-defects} has translated the familiar statement
\begin{equation*}
\partial_\mu j^\mu = 0
\end{equation*}
into the defect-language statement that the symmetry is implemented by a deformable codimension-$1$ operator whose crossing with a point insertion produces the group action. Section~8.2 will repeat the same logic for a general $q$-form symmetry, where the charged operators are extended objects rather than local fields.
\end{remark}

\section{\texorpdfstring{$q$}{q}-form symmetries and the basic dictionary}\label{sec:qform-dictionary}

The transition from ordinary symmetries to higher-form symmetries is conceptually straightforward once the codimension-$1$ defect viewpoint has been adopted. One replaces pointlike charged operators by extended charged operators and keeps the same topological-deformation logic. The only bookkeeping change is dimensional: a $q$-form symmetry acts on $q$-dimensional operators and is implemented by a topological defect of codimension $q+1$ \cite{GKSW2015,Bhardwaj2023Lectures}.

\subsection{General definition}

\begin{assumption}[Standing assumptions for Section~\ref{sec:qform-dictionary}]\label{ass:qform-standing}
Let $X$ be an oriented $d$-dimensional spacetime manifold. For the continuous abelian case, assume the theory has a conserved $(q+1)$-form current
\begin{equation*}
J^{(q+1)}\in \Omega^{q+1}(X)
\end{equation*}
satisfying
\begin{equation}\label{eq:qform-conservation}
d(\star J^{(q+1)})=0.
\end{equation}
The Hodge dual $\star J^{(q+1)}$ is therefore a closed $(d-q-1)$-form.
\end{assumption}

\begin{definition}[Continuous abelian \texorpdfstring{$q$}{q}-form symmetry operator]\label{def:qform-defect}
Let $M^{d-q-1}\subset X$ be a closed oriented codimension-$(q+1)$ submanifold. For a parameter $g=e^{i\alpha}$ in the chosen one-parameter subgroup, define
\begin{equation}\label{eq:Uq-general}
U_g(M^{d-q-1})
:=
\exp\!\left(i\alpha\int_{M^{d-q-1}}\star J^{(q+1)}\right).
\end{equation}
This is the topological defect operator implementing the continuous $q$-form symmetry.
\end{definition}

\begin{remark}[Dimensional bookkeeping]\label{rem:qform-dimensions}
The dimensions are worth recording once and for all.
\begin{enumerate}
\item The conserved current has form degree $q+1$.
\item Its Hodge dual has degree $d-q-1$ and can therefore be integrated over a manifold of dimension $d-q-1$.
\item The symmetry operator is thus supported on codimension $q+1$.
\item The charged operator has support dimension $q$ in spacetime.
\end{enumerate}
This is the higher-form generalization of the $q=0$ story from Section~\ref{sec:zero-form-defects}.
\end{remark}

\begin{proposition}[Topological invariance for a continuous \texorpdfstring{$q$}{q}-form symmetry]\label{prop:qform-topological}
If $M_0^{d-q-1}$ and $M_1^{d-q-1}$ are homologous inside a region that contains no charged $q$-dimensional operator insertions, then
\begin{equation}\label{eq:qform-invariance}
U_g(M_0^{d-q-1})=U_g(M_1^{d-q-1}).
\end{equation}
\end{proposition}

\begin{proof}
The proof is the same Stokes-theorem argument as in Proposition~\ref{prop:zero-form-topological}. If
\begin{equation*}
\partial W=\overline{M_0^{d-q-1}}\sqcup M_1^{d-q-1},
\end{equation*}
then
\begin{align*}
\int_{M_1^{d-q-1}}\star J^{(q+1)}-\int_{M_0^{d-q-1}}\star J^{(q+1)}
&=
\int_{\partial W}\star J^{(q+1)} \\
&=
\int_W d(\star J^{(q+1)}) \\
&= 0
\end{align*}
by \eqref{eq:qform-conservation}. Exponentiating proves \eqref{eq:qform-invariance}.
\end{proof}

\begin{definition}[Charged \texorpdfstring{$q$}{q}-dimensional operator]\label{def:qdim-operator}
A charged operator for a $q$-form symmetry is an operator supported on a closed oriented $q$-dimensional submanifold
\begin{equation*}
N^q\subset X.
\end{equation*}
The defining property is that when the topological defect $U_g(M^{d-q-1})$ is dragged across $N^q$, the correlation function picks up the corresponding group action.
\end{definition}

\noindent In the continuous abelian case this crossing relation is the direct analogue of \eqref{eq:crossing-relation}. If $M_{<}^{d-q-1}$ and $M_{>}^{d-q-1}$ differ by crossing a single charged operator $\mathcal W_Q(N^q)$ of charge $Q$, then
\begin{equation}\label{eq:qform-crossing}
U_g(M_{>}^{d-q-1})\,\mathcal W_Q(N^q)
=
e^{i\alpha Q}\,\mathcal W_Q(N^q)\,U_g(M_{<}^{d-q-1}).
\end{equation}
The only difference from the $q=0$ case is that the charged insertion is no longer pointlike.

\subsection{Dictionary table}

Table~\ref{tab:qform-dictionary} summarizes the pattern that will be used repeatedly in the rest of the paper.

\begin{table}[h]
    \centering
    \small
    \begin{tabular}{@{}p{0.23\linewidth}p{0.10\linewidth}p{0.20\linewidth}p{0.12\linewidth}p{0.15\linewidth}p{0.12\linewidth}@{}}
        \toprule
        Charged object & Dimension & Symmetry defect & Codim of defect & Conserved current & Background field \\
        \midrule
        Local operator $\phi(x)$ & $0$ & $U_g(\Sigma^{d-1})$ & $1$ & $J^{(1)}$ & $A^{(1)}$ \\
        Wilson line $W(C)$ & $1$ & $U_g(\Sigma^{d-2})$ & $2$ & $J^{(2)}$ & $B^{(2)}$ \\
        Surface operator $\mathcal S(\Sigma)$ & $2$ & $U_g(\Sigma^{d-3})$ & $3$ & $J^{(3)}$ & $C^{(3)}$ \\
        \bottomrule
    \end{tabular}
    \caption{Basic dictionary for continuous abelian $q$-form symmetries. The charged operator has support dimension $q$, while the symmetry defect has codimension $q+1$.}
    \label{tab:qform-dictionary}
\end{table}

\begin{remark}[Background fields]\label{rem:qform-background}
The last column of Table~\ref{tab:qform-dictionary} records the background gauge field appropriate for coupling to the symmetry. A continuous abelian $q$-form symmetry couples to a background $(q+1)$-form field, schematically through a term of the form
\begin{equation}\label{eq:qform-background-coupling}
\int_X B^{(q+1)}\wedge \star J^{(q+1)}.
\end{equation}
For discrete higher-form symmetries the background field is no longer an ordinary differential form globally, but the degree remains the same. This is the viewpoint that Chapter~\ref{ch:dw} already foreshadowed through discrete $\Z_N$ gauge fields and that Section~8.3 will make more explicit in the gauging discussion.
\end{remark}

\subsection{Example: Maxwell theory}

The cleanest higher-form example is source-free Maxwell theory in four dimensions. Let $F=dA$ be the field strength of a compact $U(1)$ gauge field. Away from electrically and magnetically charged insertions,
\begin{equation}\label{eq:maxwell-closed}
dF=0,
\qquad
d(\star F)=0.
\end{equation}
Therefore the theory has two conserved $2$-form currents:
\begin{equation}\label{eq:maxwell-currents}
J_{\mathrm{mag}}^{(2)}:=\frac{1}{2\pi}F,
\qquad
J_{\mathrm{elec}}^{(2)}:=\frac{1}{2\pi}\star F.
\end{equation}
These define a magnetic and an electric $U(1)$ $1$-form symmetry, respectively \cite{GKSW2015,Bhardwaj2023Lectures}.

\begin{definition}[Maxwell electric and magnetic symmetry operators]\label{def:maxwell-one-form}
For a closed oriented surface $\Sigma^2\subset X^4$, define
\begin{equation}\label{eq:maxwell-defects}
U^{\mathrm{mag}}_{\eta}(\Sigma)
:=
\exp\!\left(\frac{i\eta}{2\pi}\int_\Sigma F\right),
\qquad
U^{\mathrm{elec}}_{\alpha}(\Sigma)
:=
\exp\!\left(\frac{i\alpha}{2\pi}\int_\Sigma \star F\right).
\end{equation}
The first is the magnetic $1$-form symmetry defect, and the second is the electric one.
\end{definition}

\begin{remark}[What they act on]\label{rem:maxwell-acts-on-lines}
The electric $1$-form symmetry acts on Wilson lines
\begin{equation*}
W_n(C)=\exp\!\left(i n\oint_C A\right),
\end{equation*}
while the magnetic $1$-form symmetry acts on 't~Hooft lines. In both cases the charged operator has support dimension $1$, exactly as predicted by Table~\ref{tab:qform-dictionary}.
\end{remark}

\begin{proposition}[Maxwell realizes both electric and magnetic \texorpdfstring{$1$}{1}-form symmetries]\label{prop:maxwell-one-form}
In source-free compact Maxwell theory in four dimensions, both operators in \eqref{eq:maxwell-defects} are topological away from charged line insertions. When the electric defect $U^{\mathrm{elec}}_{\alpha}(\Sigma)$ is dragged across a Wilson line of charge $n$, the correlator picks up the phase $e^{i\alpha n}$. When the magnetic defect $U^{\mathrm{mag}}_{\eta}(\Sigma)$ is dragged across an 't~Hooft line of magnetic charge $m$, the correlator picks up $e^{i\eta m}$.
\end{proposition}

\begin{proof}
Topological invariance follows immediately from Proposition~\ref{prop:qform-topological} applied to the currents \eqref{eq:maxwell-currents} and the conservation laws \eqref{eq:maxwell-closed}.

\medskip

\noindent For the electric symmetry, let $\Sigma_{<}$ and $\Sigma_{>}$ be two closed surfaces that differ by moving the defect across a Wilson line $W_n(C)$. Choose a small $2$-sphere $S^2_C$ linking $C$ once and a $3$-manifold $W$ with boundary
\begin{equation*}
\partial W=\overline{\Sigma_{<}}\sqcup \Sigma_{>}\sqcup \overline{S^2_C}.
\end{equation*}
Applying Stokes' theorem away from the line insertion gives
\begin{align*}
0
&=
\int_W d(\star F) \\
&=
\int_{\Sigma_{>}}\star F
-\int_{\Sigma_{<}}\star F
-\int_{S^2_C}\star F.
\end{align*}
Hence
\begin{equation*}
\frac{1}{2\pi}\int_{\Sigma_{>}}\star F
=
\frac{1}{2\pi}\int_{\Sigma_{<}}\star F
+
\frac{1}{2\pi}\int_{S^2_C}\star F.
\end{equation*}
For a Wilson line of electric charge $n$, the last term is $n$. Therefore
\begin{equation*}
U^{\mathrm{elec}}_{\alpha}(\Sigma_{>})\,W_n(C)
=
e^{i\alpha n}\,W_n(C)\,U^{\mathrm{elec}}_{\alpha}(\Sigma_{<}).
\end{equation*}

\medskip

\noindent The magnetic case is identical with $F$ in place of $\star F$. If $S^2_{C_m}$ links an 't~Hooft line of magnetic charge $m$ once, then
\begin{equation*}
\frac{1}{2\pi}\int_{S^2_{C_m}} F = m,
\end{equation*}
so dragging $U^{\mathrm{mag}}_{\eta}(\Sigma)$ across that line produces the phase $e^{i\eta m}$.
\end{proof}

\subsection{Example: the center symmetry of pure Yang--Mills}

The Maxwell example is continuous and abelian. Pure non-abelian Yang--Mills supplies the other important example for this paper: a discrete $1$-form symmetry. Consider a pure $SU(N)$ gauge theory in four dimensions, with no matter fields in representations of nonzero $N$-ality. Then every local field transforms in the adjoint representation and is therefore neutral under the center $\Z_N\subset SU(N)$. The center does, however, act nontrivially on Wilson lines. This is exactly the hallmark of a $1$-form symmetry \cite{GKSW2015,Bhardwaj2023Lectures}.

\begin{example}[Center symmetry as a \texorpdfstring{$1$}{1}-form symmetry]\label{ex:center-symmetry}
Let
\begin{equation*}
W_R(C)
\end{equation*}
be a Wilson line in a representation $R$ of $SU(N)$, and let $k_R\in \Z_N$ denote its $N$-ality. The discrete center symmetry is implemented by a codimension-$2$ topological operator
\begin{equation*}
U_z(\Sigma^{d-2}),
\qquad
z\in \Z_N,
\end{equation*}
whose crossing with the Wilson line is
\begin{equation}\label{eq:center-crossing}
U_z(\Sigma_{>})\,W_R(C)
=
z^{k_R}\,W_R(C)\,U_z(\Sigma_{<}).
\end{equation}
Thus the Wilson line is the charged object, the symmetry defect is a surface operator, and the symmetry is a discrete $1$-form symmetry rather than an ordinary $0$-form symmetry.
\end{example}

\begin{remark}[Why local operators are neutral]\label{rem:center-neutral}
Every local gauge-invariant operator in pure Yang--Mills is built from adjoint fields such as $F_{\mu\nu}$ and covariant derivatives thereof. Since adjoint fields have zero $N$-ality, no local operator carries the center charge. This is why the center symmetry cannot be an ordinary global symmetry acting on point operators. Its charged objects are instead Wilson lines.
\end{remark}

\begin{remark}[Forward pointer to gauging]\label{rem:qform-gauging-forward}
Section~\ref{sec:qform-dictionary} has provided the dictionary we will need later: support dimension of charged operators, codimension of symmetry defects, current degree, and background-field degree. Section~8.3 will use exactly this dictionary to explain what it means to gauge a higher-form symmetry, first abstractly and then in the concrete $\Z_N$ examples related to Chapters~\ref{ch:bf} and \ref{ch:dw}.
\end{remark}

\section{Gauging higher-form symmetries}\label{sec:qform-gauging}

Sections~\ref{sec:zero-form-defects} and \ref{sec:qform-dictionary} treated symmetry operators as probes: one couples the theory to a background field and studies how charged operators respond. Gauging the symmetry means promoting that background data to dynamical data and summing over it in the path integral. For higher-form symmetries, the logic is the same as for ordinary symmetries; only the degree of the background field changes \cite{GKSW2015,Bhardwaj2023Lectures}.

\subsection{General recipe}

\begin{assumption}[Finite abelian \texorpdfstring{$q$}{q}-form symmetry]\label{ass:qform-gauging}
Let $X$ be a closed oriented $d$-manifold, and suppose the theory has a finite abelian $q$-form symmetry $A$. Assume that coupling to a background $A$ gauge field is already defined, so for each background configuration
\begin{equation*}
B\in H^{q+1}(X,A)
\end{equation*}
there is a partition function
\begin{equation*}
Z(X;B).
\end{equation*}
We use the same orbit-weighted convention as in Chapter~\ref{ch:dw}: backgrounds are summed over up to gauge equivalence, with a finite automorphism factor.
\end{assumption}

\begin{definition}[Gauging a finite \texorpdfstring{$q$}{q}-form symmetry]\label{def:gauging-qform}
Under Assumption~\ref{ass:qform-gauging}, the gauged theory is defined by
\begin{equation}\label{eq:qform-gauge-sum}
Z_{X/A}
:=
\sum_{[B]\in H^{q+1}(X,A)}
\frac{1}{|\Aut(B)|}\,
Z(X;B).
\end{equation}
Equivalently, one promotes the background $(q+1)$-form gauge field to a dynamical field and sums over its gauge-equivalence classes.
\end{definition}

\begin{remark}[Continuous symmetries]\label{rem:qform-continuous-gauging}
For a continuous abelian $q$-form symmetry, the discrete sum \eqref{eq:qform-gauge-sum} is replaced by a functional integral over a dynamical $(q+1)$-form gauge field, modulo its higher-form gauge redundancy. The present paper will only need the finite case, because that is the case most directly tied to the $\Z_N$ examples of Chapters~\ref{ch:bf} and \ref{ch:dw}.
\end{remark}

\begin{proposition}[Effect on charged operators]\label{prop:gauging-selection}
If an operator is charged under the symmetry being gauged, then it is not a genuine operator of the gauged theory unless it can be dressed so as to become gauge-invariant. In particular, summing over the background field projects onto the symmetry-invariant sector.
\end{proposition}

\begin{proof}
Before gauging, a charged operator $\mathcal W_Q$ transforms under the symmetry defect by a phase, as in \eqref{eq:qform-crossing}. After gauging, the background field $B$ is summed over. Gauge invariance of the full path integral requires invariance under higher-form gauge transformations of $B$. An insertion with nontrivial charge changes by a nontrivial phase under those transformations and therefore does not define a gauge-invariant operator by itself. Equivalently, the sum over $B$ averages over the symmetry orbit and projects away non-invariant insertions.
\end{proof}

\subsection{Worked example: gauging the electric \texorpdfstring{$\Z_N$}{Z\_N} subgroup in four-dimensional Maxwell theory}

We now work through the most concrete example in full. Start from four-dimensional compact Maxwell theory with gauge field $A$ and field strength $F=dA$. Section~\ref{sec:qform-dictionary} already identified the electric $U(1)^{(1)}$ symmetry acting on Wilson lines. We now gauge the finite subgroup
\begin{equation*}
\Z_N \subset U(1)^{(1)}_{\mathrm{elec}}.
\end{equation*}

\begin{assumption}[Background field for the electric \texorpdfstring{$\Z_N$}{Z\_N} one-form symmetry]\label{ass:maxwell-zn-gauging}
Let $X$ be a closed oriented four-manifold. A background field for the electric $\Z_N$ one-form symmetry is represented locally by a two-form $B$ with periods in
\begin{equation*}
\frac{2\pi}{N}\Z.
\end{equation*}
It is accompanied by a one-form gauge redundancy
\begin{equation}\label{eq:one-form-gauge-transformation}
A \longmapsto A+\lambda,
\qquad
B \longmapsto B+d\lambda,
\end{equation}
where $\lambda$ is a $U(1)$ connection whose allowed periods are integer multiples of $2\pi/N$.
\end{assumption}

\begin{definition}[Maxwell theory coupled to the electric background]\label{def:maxwell-background}
One convenient background-coupled action is
\begin{equation}\label{eq:maxwell-B-action}
S[A,B]
:=
\frac{1}{2e^2}\int_X (F-B)\wedge \star(F-B).
\end{equation}
Because the combination $F-B$ is invariant under \eqref{eq:one-form-gauge-transformation}, the action depends only on the $\Z_N$ background gauge field and not on the auxiliary local choice of $B$ and $\lambda$ separately.
\end{definition}

\begin{proposition}[Gauging projects Wilson lines onto multiples of \texorpdfstring{$N$}{N}]\label{prop:maxwell-zn-selection}
After gauging the electric $\Z_N$ one-form symmetry of four-dimensional compact Maxwell theory, a Wilson line
\begin{equation}\label{eq:maxwell-wilson-n}
W_n(C)=\exp\!\left(i n\oint_C A\right)
\end{equation}
is a genuine line operator only if
\begin{equation}\label{eq:n-multiple-N}
n\equiv 0 \pmod N.
\end{equation}
\end{proposition}

\begin{proof}
Under the one-form gauge transformation \eqref{eq:one-form-gauge-transformation}, the Wilson line transforms as
\begin{equation}\label{eq:wilson-one-form-transform}
W_n(C)\longmapsto
\exp\!\left(i n\oint_C \lambda\right)\,W_n(C).
\end{equation}
Since $\lambda$ is allowed to have periods in $(2\pi/N)\Z$, there are allowed gauge transformations for which
\begin{equation*}
\oint_C \lambda = \frac{2\pi}{N}.
\end{equation*}
For \eqref{eq:wilson-one-form-transform} to be invariant under every allowed transformation, we must have
\begin{equation*}
\exp\!\left(\frac{2\pi i n}{N}\right)=1,
\end{equation*}
which is equivalent to \eqref{eq:n-multiple-N}. Thus only the lines $W_{N\ell}(C)$ survive as genuine operators in the gauged theory.
\end{proof}

\begin{remark}[What changes, and what does not]\label{rem:maxwell-global-form}
Equation \eqref{eq:n-multiple-N} is often summarized by saying that gauging the electric $\Z_N$ one-form symmetry changes the global form from $U(1)$ to $U(1)/\Z_N$. Since these groups are isomorphic as Lie groups, the local gauge algebra is unchanged. What changes is the global structure encoded by the spectrum of genuine line operators. That is the physically meaningful statement in the present context.
\end{remark}

\begin{remark}[The dual symmetry after gauging]\label{rem:maxwell-dual-symmetry}
In $d$ spacetime dimensions, gauging a finite $q$-form symmetry produces a dual finite $(d-q-2)$-form symmetry. For the present example,
\begin{equation*}
d=4,
\qquad
q=1,
\end{equation*}
so the dual symmetry is again a $1$-form symmetry, not a $0$-form symmetry. Its charged objects are magnetic line operators. The common slogan about pointlike monopole operators applies after reducing to three dimensions, but not to the four-dimensional theory considered here.
\end{remark}

\subsection{Relation to \texorpdfstring{$SU(N)$}{SU(N)} and \texorpdfstring{$PSU(N)$}{PSU(N)} gauge theories}

The same operation appears in a non-abelian setting. Pure $SU(N)$ Yang--Mills theory has an electric center symmetry
\begin{equation*}
\Z_N^{(1)}
\end{equation*}
acting on Wilson lines by their $N$-ality, as explained in Example~\ref{ex:center-symmetry}. Gauging this one-form symmetry removes Wilson lines of nonzero $N$-ality from the set of genuine line operators and replaces the global form of the gauge group by
\begin{equation*}
PSU(N)=SU(N)/\Z_N.
\end{equation*}

\begin{remark}[Discrete theta-angle caveat]\label{rem:psun-discrete-theta}
This statement needs one qualification. Gauging the center symmetry of pure $SU(N)$ Yang--Mills does not produce a unique theory; in general one must also choose a discrete theta angle. The simplest choice is the untwisted case, often denoted $PSU(N)$ with discrete theta parameter $p=0$. More generally one obtains a family of $PSU(N)_p$ theories labeled by
\begin{equation*}
p\in \Z_N
\end{equation*}
\cite{AharonySeibergTachikawa2013}. For the purposes of the present review, the important point is the structural one: gauging the electric center one-form symmetry changes the global form and the genuine line spectrum.
\end{remark}

\begin{remark}[Forward pointer to Dijkgraaf--Witten and BF theory]\label{rem:qform-to-dw}
The finite sum \eqref{eq:qform-gauge-sum} is the same structural move that appeared in Chapter~\ref{ch:dw}: one sums over discrete gauge-field backgrounds with the appropriate gauge-volume factor. If the infrared dynamics is fully gapped and only the finite higher-form gauge field survives, the resulting theory is naturally described by a discrete TQFT of the sort studied in Chapters~\ref{ch:bf} and \ref{ch:dw}. Section~8.4 will use this viewpoint again when we define anomalies and then reinterpret twisted Dijkgraaf--Witten theory as a clean inflow example.
\end{remark}

\section{Anomalies as obstructions to gauging}\label{sec:anomaly-definition}

The previous section explained how a symmetry is gauged when the background-field dependence is honest gauge-invariant data. An anomaly is the obstruction to doing that. Operationally, one couples the theory to a background field for the symmetry, performs a background gauge transformation, and asks whether the partition function is unchanged. If it is not unchanged, but instead shifts by a local functional of the background field, then the symmetry has an 't~Hooft anomaly \cite{GKSW2015,Bhardwaj2023Lectures}.

\subsection{Operational definition}

\begin{assumption}[Background-field setup]\label{ass:anomaly-background}
Let $X$ be a closed oriented spacetime manifold, let $G$ be a global symmetry of the theory, and let $A$ denote a background gauge field for $G$ in the sense appropriate to that symmetry. Write
\begin{equation*}
Z[X;A]
\end{equation*}
for the partition function in the presence of that background.
\end{assumption}

\begin{definition}['t~Hooft anomaly]\label{def:tHooft-anomaly}
The theory has an \emph{'t~Hooft anomaly} for the symmetry $G$ if, under a background gauge transformation
\begin{equation*}
A\longmapsto A^g,
\end{equation*}
the partition function transforms as
\begin{equation}\label{eq:anomaly-phase}
Z[X;A^g]
=
Z[X;A]\,
\exp\!\bigl(i\alpha_X[A,g]\bigr),
\end{equation}
where $\alpha_X[A,g]$ is a local functional of the background field and the gauge-transformation parameter, and there is no local counterterm
\begin{equation*}
S_{\mathrm{ct}}[A]
\end{equation*}
for which the redefined partition function
\begin{equation*}
Z'[X;A]:=Z[X;A]\,e^{iS_{\mathrm{ct}}[A]}
\end{equation*}
becomes gauge-invariant.
\end{definition}

\begin{remark}[Why locality matters]\label{rem:anomaly-locality}
The word ``local'' in Definition~\ref{def:tHooft-anomaly} is essential. Multiplying the partition function by an arbitrary nonlocal phase would be physically meaningless bookkeeping. The anomaly data is the part that survives after allowing every possible \emph{local} counterterm built from the background field. This is why anomalies are robust and why they can be matched across renormalization-group flow.
\end{remark}

\begin{proposition}[An anomaly obstructs gauging]\label{prop:anomaly-obstruction}
If the symmetry $G$ has an 't~Hooft anomaly in the sense of Definition~\ref{def:tHooft-anomaly}, then the naive gauging prescription fails: one cannot consistently sum over background-field configurations to produce a gauge-invariant theory of a dynamical gauge field for $G$.
\end{proposition}

\begin{proof}
The gauging prescription of Section~\ref{sec:qform-gauging} requires summing over gauge-equivalence classes of background fields. For that sum to be well defined, the integrand must assign the same weight to gauge-equivalent configurations. But by \eqref{eq:anomaly-phase}, the weights for $A$ and $A^g$ differ by the phase
\begin{equation*}
\exp\!\bigl(i\alpha_X[A,g]\bigr).
\end{equation*}
If this phase could be removed by a local counterterm, the anomaly would be trivial by definition. If it cannot be removed, then gauge-equivalent configurations contribute different amplitudes, and the quotient by gauge equivalence is not well defined. Therefore the naive gauging procedure is obstructed.
\end{proof}

\begin{remark}[What an anomaly does \emph{not} mean]\label{rem:anomaly-not-inconsistency}
An 't~Hooft anomaly is not an inconsistency of the original ungauged theory. The theory is perfectly well defined as long as $G$ is treated as a global symmetry and the background field is kept nondynamical. The inconsistency appears only when one attempts to gauge the symmetry without supplying extra higher-dimensional inflow data.
\end{remark}

\subsection{Warmup example: two-dimensional chiral fermions}

The standard field-theory warmup is a two-dimensional massless Dirac fermion. Classically it has both vector and axial symmetries,
\begin{equation*}
U(1)_V\times U(1)_A,
\end{equation*}
but quantum mechanically one cannot gauge both simultaneously. This is the simplest operational example of the anomaly definition above and the prototype for the fuller inflow discussion later in the chapter.

\begin{example}[Simultaneous gauging is obstructed]\label{ex:2d-chiral-warmup}
Let $\psi$ be a massless Dirac fermion in two dimensions. If one couples the vector current and the axial current to background fields $A_V$ and $A_A$, then a regularization that preserves vector gauge invariance makes the axial current anomalous, while a regularization that preserves axial gauge invariance makes the vector current anomalous. In either scheme, one cannot keep both background symmetries gauge-invariant simultaneously. Equivalently, the background-field partition function acquires a local anomalous phase under one of the two gauge transformations.
\end{example}

\begin{remark}[Why this example is only a warmup]\label{rem:2d-warmup-scope}
We are not deriving the full two-dimensional anomaly polynomial here. The only point needed in the present section is the logic: the anomalous variation is local, it cannot be removed while preserving both symmetries at once, and therefore simultaneous gauging is obstructed. Section~8.6 will return to the high-energy version of this story in the more systematic anomaly-inflow framework connected to Schwartz's anomaly chapter.
\end{remark}

\begin{remark}[Forward pointer to inflow]\label{rem:anomaly-to-inflow}
Definition~\ref{def:tHooft-anomaly} treats the anomalous phase as an obstruction. The next section will convert that obstruction into a positive statement: the anomalous boundary variation can be cancelled by a higher-dimensional bulk term. Chapter~\ref{ch:dw} has already shown this mechanism in a discrete setting; Sections~8.5--8.7 will now place it in the general language of anomaly inflow.
\end{remark}

\section{Anomaly inflow: general statement}\label{sec:anomaly-inflow-general}

Section~\ref{sec:anomaly-definition} defined an anomaly as the failure of the boundary partition function to be gauge-invariant in background fields. Anomaly inflow is the statement that this failure can sometimes be cancelled by placing the anomalous theory on the boundary of a one-dimension-higher bulk whose own variation produces precisely the opposite phase. The anomaly is therefore not removed; it is reinterpreted as the boundary avatar of a bulk topological response \cite{CallanHarvey1985,Freed1995CSnotes,FreedMooreTeleman2024}.

\subsection{Setup}

\begin{assumption}[Bulk-boundary geometry]\label{ass:inflow-setup}
Let $M$ be a compact oriented $(d+1)$-manifold with boundary
\begin{equation*}
\partial M = X.
\end{equation*}
Suppose a background field $A$ for the symmetry of interest is defined on $M$, and let $A|_X$ denote its restriction to the boundary.
\end{assumption}

\begin{definition}[Invertible bulk TQFT]\label{def:invertible-bulk}
An \emph{invertible} $(d+1)$-dimensional TQFT is one whose partition function on a closed manifold is a nonzero complex phase and whose Hilbert space on every closed spatial slice is one-dimensional. Equivalently, the bulk contributes no propagating topological sectors of its own; it supplies only a phase.
\end{definition}

\begin{remark}[Why invertibility matters]\label{rem:invertible-why}
If the bulk were not invertible, then the combined bulk-boundary system would contain genuine bulk topological degrees of freedom, not merely a bookkeeping device for the boundary anomaly. That is sometimes physically interesting, but it is not the minimal inflow mechanism. For anomaly cancellation, the cleanest setup is therefore an invertible bulk whose only job is to contribute the compensating phase.
\end{remark}

\noindent Let
\begin{equation*}
Z_{\partial}[X;A|_X]
\end{equation*}
be the partition function of the anomalous $d$-dimensional theory on the boundary, and let
\begin{equation*}
Z_{\mathrm{bulk}}[M;A]
=
\exp\!\bigl(iS_{\mathrm{inv}}[M;A]\bigr)
\end{equation*}
be the partition function of the invertible bulk theory.

\subsection{Cancellation of the anomalous variation}

\begin{proposition}[General inflow mechanism]\label{prop:general-inflow}
Assume the boundary theory has anomaly phase
\begin{equation}\label{eq:boundary-anomaly-phase}
Z_{\partial}[X;(A|_X)^g]
=
Z_{\partial}[X;A|_X]\,
\exp\!\bigl(i\alpha_X[A|_X,g]\bigr),
\end{equation}
and assume the bulk partition function on a manifold with boundary transforms as
\begin{equation}\label{eq:bulk-anomaly-phase}
Z_{\mathrm{bulk}}[M;A^g]
=
Z_{\mathrm{bulk}}[M;A]\,
\exp\!\bigl(-i\alpha_X[A|_X,g]\bigr).
\end{equation}
Then the combined bulk-boundary system
\begin{equation}\label{eq:combined-partition}
Z_{\mathrm{tot}}[M,X;A]
:=
Z_{\mathrm{bulk}}[M;A]\,
Z_{\partial}[X;A|_X]
\end{equation}
is gauge-invariant.
\end{proposition}

\begin{proof}
Apply the background gauge transformation to \eqref{eq:combined-partition}. Using \eqref{eq:boundary-anomaly-phase} and \eqref{eq:bulk-anomaly-phase}, we find
\begin{align*}
Z_{\mathrm{tot}}[M,X;A^g]
&=
Z_{\mathrm{bulk}}[M;A^g]\,
Z_{\partial}[X;(A|_X)^g] \\
&=
Z_{\mathrm{bulk}}[M;A]\,
e^{-i\alpha_X[A|_X,g]}\,
Z_{\partial}[X;A|_X]\,
e^{i\alpha_X[A|_X,g]} \\
&=
Z_{\mathrm{bulk}}[M;A]\,
Z_{\partial}[X;A|_X] \\
&=
Z_{\mathrm{tot}}[M,X;A].
\end{align*}
Thus the anomalous boundary variation is cancelled by the bulk variation.
\end{proof}

\begin{remark}[What the bulk variation means]\label{rem:bulk-closed-vs-boundary}
There is no contradiction in calling the bulk theory gauge-invariant on closed manifolds and yet gauge-variant on manifolds with boundary. On a closed manifold, the anomalous phase in \eqref{eq:bulk-anomaly-phase} has nowhere to live because there is no boundary local functional. The variation appears only when a boundary is present, and then it is precisely a boundary term.
\end{remark}

\begin{remark}[Relation to Chapter~\ref{ch:dw}]\label{rem:inflow-dw-relation}
Section~\ref{sec:dw-inflow} already exhibited the same logic in a discrete finite-group setting. There, a nontrivial Dijkgraaf--Witten cocycle obstructed a standalone gauge-invariant boundary completion, and the bulk cocycle supplied the compensating variation. Proposition~\ref{prop:general-inflow} is the continuum-language version of that same structure.
\end{remark}

\subsection{The standard high-energy template}

The standard high-energy example is a four-dimensional chiral anomaly cancelled by a five-dimensional bulk Chern--Simons-type term. Concretely, one considers a five-dimensional action built from the background gauge field and its field strength, schematically of the form
\begin{equation}\label{eq:five-d-cs-schematic}
\int_M A\wedge F\wedge F,
\end{equation}
or its nonabelian generalization. On a manifold with boundary, the gauge variation of \eqref{eq:five-d-cs-schematic} is a four-dimensional boundary term. Section~8.6 will derive that variation explicitly and match it to the four-dimensional chiral anomaly.

\begin{remark}[Condensed-matter counterpart]\label{rem:cm-inflow-forward}
The same mechanism appears in condensed matter, with one dimension lower: the fractional quantum Hall edge is anomalous by itself, but its anomalous variation is cancelled by the bulk Chern--Simons response term. Section~8.7 will treat that case after the high-energy derivation, so that the two examples can be compared directly.
\end{remark}

\section{High-energy example: the chiral anomaly as five-dimensional inflow}\label{sec:hep-inflow}

We now connect the general inflow statement to the anomaly material from Schwartz's Chapter~30. The cleanest direct match is the axial anomaly of a four-dimensional Dirac fermion in an external gauge field. In that case the inflow bulk is most naturally written as a \emph{mixed} five-dimensional Chern--Simons-type term involving the axial background field and the vector field strength. This is slightly narrower than the full nonabelian $5$-dimensional Chern--Simons form, but it is the version that matches Schwartz's equation most directly and with the least convention drift.

\subsection{Boundary anomaly in Schwartz's notation}

\begin{assumption}[Boundary theory and conventions]\label{ass:hep-inflow}
Let $X$ be a closed oriented four-manifold bounding a compact oriented five-manifold $N$:
\begin{equation*}
\partial N = X.
\end{equation*}
Let $A$ be a background vector gauge field on $N$ with field strength
\begin{equation*}
F=dA+A\wedge A,
\end{equation*}
and let $a$ be a background axial $U(1)$ gauge field on $N$, restricting on the boundary to the source for the axial current $j_A^\mu$.
\end{assumption}

\noindent In the component convention quoted in Task~8.6 and in Schwartz Chapter~30, the axial anomaly is
\begin{equation}\label{eq:schwartz-3020-task}
\partial_\mu j_A^\mu
=
\frac{1}{16\pi^2}\,
\varepsilon^{\mu\nu\rho\sigma}\,
\tr(F_{\mu\nu}F_{\rho\sigma}).
\end{equation}
Using the four-dimensional identity
\begin{equation}\label{eq:FF-component-form}
\tr(F\wedge F)
=
\frac{1}{4}\,
\varepsilon^{\mu\nu\rho\sigma}\,
\tr(F_{\mu\nu}F_{\rho\sigma})\,d^4x,
\end{equation}
equation \eqref{eq:schwartz-3020-task} is equivalently
\begin{equation}\label{eq:axial-anomaly-forms}
d(\star j_A)=\frac{1}{4\pi^2}\,\tr(F\wedge F).
\end{equation}

\begin{remark}[What is being matched]\label{rem:schwartz-match}
Equation \eqref{eq:axial-anomaly-forms} is the precise form that inflow must reproduce. The coefficient is fixed entirely by \eqref{eq:schwartz-3020-task} together with the form identity \eqref{eq:FF-component-form}. This is the reason for using the mixed bulk term below: it matches the axial-current anomaly without first passing through the fuller descent formalism for the consistent nonabelian gauge anomaly.
\end{remark}

\subsection{The five-dimensional bulk term}

\begin{definition}[Mixed five-dimensional inflow action]\label{def:mixed-5d-inflow}
Define the bulk action
\begin{equation}\label{eq:mixed-5d-action}
S_{\mathrm{bulk}}[a,A]
:=
-\frac{1}{4\pi^2}\int_N a\wedge \tr(F\wedge F).
\end{equation}
\end{definition}

\begin{remark}[Why this is Chern--Simons type]\label{rem:mixed-cs-type}
The integrand in \eqref{eq:mixed-5d-action} is a $5$-form built from one gauge potential and two field strengths. It is therefore of the same Chern--Simons type as the schematic expression from Section~\ref{sec:anomaly-inflow-general},
\begin{equation*}
\int_N A\wedge F\wedge F.
\end{equation*}
For the present axial-anomaly application, the one-form potential is the axial background field $a$, while the two field strengths come from the vector gauge field.
\end{remark}

\begin{proposition}[Gauge variation of the mixed bulk term]\label{prop:mixed-bulk-variation}
Under an axial background gauge transformation
\begin{equation}\label{eq:axial-gauge-transform}
a\longmapsto a+d\lambda,
\end{equation}
with the vector background field $A$ held fixed, the bulk action varies by the boundary term
\begin{equation}\label{eq:mixed-bulk-variation}
\delta_\lambda S_{\mathrm{bulk}}
=
-\frac{1}{4\pi^2}\int_X \lambda\,\tr(F\wedge F).
\end{equation}
\end{proposition}

\begin{proof}
Since $A$ is fixed, the only variation comes from $a$:
\begin{equation*}
\delta_\lambda S_{\mathrm{bulk}}
=
-\frac{1}{4\pi^2}\int_N d\lambda\wedge \tr(F\wedge F).
\end{equation*}
Now use the Bianchi identity
\begin{equation*}
d_A F = 0,
\end{equation*}
which implies
\begin{equation*}
d\,\tr(F\wedge F)=0.
\end{equation*}
Therefore
\begin{align*}
\delta_\lambda S_{\mathrm{bulk}}
&=
-\frac{1}{4\pi^2}\int_N d\!\left(\lambda\,\tr(F\wedge F)\right) \\
&=
-\frac{1}{4\pi^2}\int_X \lambda\,\tr(F\wedge F),
\end{align*}
where the second line is Stokes' theorem with the boundary orientation induced from $N$. This is \eqref{eq:mixed-bulk-variation}.
\end{proof}

\subsection{Matching to the boundary fermion}

\begin{proposition}[The bulk variation cancels the boundary axial anomaly]\label{prop:axial-inflow-match}
Let the boundary theory be a four-dimensional fermionic theory whose axial anomaly is \eqref{eq:axial-anomaly-forms}. Then the anomalous variation of the boundary effective action under
\begin{equation*}
a|_X \longmapsto a|_X + d\lambda
\end{equation*}
is
\begin{equation}\label{eq:boundary-axial-variation}
\delta_\lambda W_{\partial}
=
\frac{1}{4\pi^2}\int_X \lambda\,\tr(F\wedge F),
\end{equation}
and this is cancelled exactly by \eqref{eq:mixed-bulk-variation}.
\end{proposition}

\begin{proof}
The background field $a|_X$ couples to the boundary axial current, so under an infinitesimal axial gauge transformation with parameter $\lambda$ the change in the boundary effective action is
\begin{equation*}
\delta_\lambda W_{\partial}
=
\int_X \lambda\,d(\star j_A).
\end{equation*}
Substituting the anomaly equation \eqref{eq:axial-anomaly-forms} gives
\begin{equation*}
\delta_\lambda W_{\partial}
=
\frac{1}{4\pi^2}\int_X \lambda\,\tr(F\wedge F),
\end{equation*}
which is exactly \eqref{eq:boundary-axial-variation}. Comparing with Proposition~\ref{prop:mixed-bulk-variation}, we see that
\begin{equation*}
\delta_\lambda W_{\partial}+\delta_\lambda S_{\mathrm{bulk}}=0.
\end{equation*}
Hence the combined bulk-boundary partition function is gauge-invariant.
\end{proof}

\begin{remark}[Callan--Harvey mechanism in the present notation]\label{rem:callan-harvey-present}
Propositions~\ref{prop:mixed-bulk-variation} and \ref{prop:axial-inflow-match} are the Callan--Harvey inflow mechanism in the present notation: the anomalous divergence of the boundary current is not an isolated failure of gauge invariance, but the boundary shadow of a higher-dimensional topological response term.
\end{remark}

\subsection{Relation to the full \texorpdfstring{$5$}{5}-dimensional Chern--Simons form}

\begin{remark}[The full nonabelian descendant story]\label{rem:full-cs5-story}
If one studies a genuinely chiral four-dimensional gauge theory, rather than the axial current of a Dirac fermion in an external vector field, the relevant bulk term is the full nonabelian Chern--Simons $5$-form
\begin{equation}\label{eq:cs5-full}
\omega_5(A)
:=
\tr\!\left(
A\wedge dA\wedge dA
+\frac{3}{2}A\wedge A\wedge A\wedge dA
+\frac{3}{5}A\wedge A\wedge A\wedge A\wedge A
\right),
\end{equation}
which satisfies
\begin{equation*}
d\omega_5(A)=\tr(F^3).
\end{equation*}
Its gauge variation on a manifold with boundary is again a boundary term, now expressed through the standard descent form. We do not derive that full descent formalism here. For the present paper, the mixed term \eqref{eq:mixed-5d-action} is the right choice because it matches Schwartz's axial-anomaly equation directly.
\end{remark}

\begin{remark}[Forward pointer]\label{rem:hep-to-cm-forward}
Section~\ref{sec:hep-inflow} is the main high-energy inflow example for the paper. Section~8.7 will now present the condensed-matter analogue, where a boundary chiral boson anomaly is cancelled by the bulk Chern--Simons response of the fractional quantum Hall fluid.
\end{remark}
